\documentclass{article}
\usepackage{PRIMEarxiv} % Core package for the PrimeAI Template
\usepackage{amsmath, amssymb, amsthm, bm, dcolumn} % Add amsthm here for the proof environment
\usepackage[numbers,sort&compress]{natbib} % Natbib for citations
\usepackage{graphicx} % For high-quality images
\usepackage{multicol} % Multi-column figures/tables
\usepackage{hyperref} % Hyperlinks
\usepackage{listings} % Code listings
\usepackage{authblk} % For structured affiliations
% Define theorem style (optional)
\theoremstyle{plain}
\newtheorem{theorem}{Theorem}
\renewcommand{\qedsymbol}{}

% Custom commands for primes and qubit encoding
\newcommand{\primeQubit}[1]{|p_{#1}\rangle}
\newcommand{\primeGate}[1]{U_{p_{#1}}}

\usepackage{wrapfig}
\usepackage[pscoord]{eso-pic}
\usepackage[fulladjust]{marginnote}
\reversemarginpar

% Typesetting improvements without footnote patching
\usepackage[protrusion=true, expansion=true, tracking=false]{microtype}
\microtypecontext{spacing=nonfrench}

% Line numbers
\usepackage[right]{lineno}

% Text layout - adjust as needed
\raggedright
\setlength{\parindent}{0.5cm}
\textwidth 5.25in 
\textheight 8.75in

% Set double spacing
\usepackage{setspace} 
\doublespacing

% Adjust width for specific content
\usepackage{changepage}

% Adjust caption style
\usepackage[aboveskip=1pt,labelfont=bf,labelsep=period,singlelinecheck=off]{caption}

% Remove brackets from references
\makeatletter
\renewcommand{\@biblabel}[1]{\quad#1.}
\makeatother

% Header, footer, and page numbers
\usepackage{lastpage,fancyhdr}
\pagestyle{fancy}
\fancyhf{}
\fancyhead[L]{Preprint - PrimeAI Enhanced Template}
\fancyfoot[C]{\scriptsize Multiplicity Theory © 2024 Ryan Van Gelder - Citizen Gardens \\ Licensed Under MIT and CC BY-NC-SA 4.0.}
\fancyfoot[R]{Page \thepage\ of \pageref{LastPage}}
\renewcommand{\footrule}{\hrule height 2pt \vspace{2mm}}

\begin{document}

\title{Helium Bose Einstein Condensate\\With Multiplicity Theory}

\author{Ryan O. Van Gelder}
\affil{Citizen Gardens - The Foundation of Multiplicity \\ \texttt{ryan@citizengardens.org}}
\date{\today}

\maketitle

\begin{abstract}
This paper explores the integration of Helium Bose-Einstein Condensates (BEC) within the framework of Multiplicity Theory. By examining the foundational principles of quantum condensates \cite{bose1924,einstein1925} and their alignment with prime-encoded structures, this study develops recursive feedback models and evaluates stability conditions for superfluid dynamics. Applications in quantum computing and astrophysical simulations are discussed, leveraging the coherence and macroscopic quantum phenomena of Helium BECs \cite{anderson1995,greiner2002}.
\end{abstract}

\section{Introduction}
The Bose-Einstein Condensate (BEC), first predicted by Satyendra Nath Bose and Albert Einstein \cite{bose1924,einstein1925}, is a state of matter formed when bosons occupy the same quantum ground state at near absolute zero temperatures. Experimental realization in 1995 by Anderson et al. confirmed this quantum phenomenon, marking a milestone in condensed matter physics \cite{anderson1995}.

Helium BECs exhibit superfluidity and coherence, making them ideal candidates for integration with Multiplicity Theory. This framework encodes interactions using prime numbers to model recursive systems \cite{multiplicity2024}. This study investigates:
\begin{itemize}
    \item Coherence alignment with prime-based eigenstates.
    \item Recursive interaction matrices for dynamic modeling.
    \item Applications in quantum simulations and astrophysical systems.
\end{itemize}

\section{Theoretical Foundations}
\subsection{Bose-Einstein Condensation}
A BEC occurs when particles obey Bose-Einstein statistics, allowing multiple bosons to occupy the same quantum state. The condensate wavefunction $\psi$ represents the macroscopic quantum state:
\begin{equation}
\psi(\mathbf{r},t) = \sqrt{n(\mathbf{r},t)} e^{i \phi(\mathbf{r},t)},
\end{equation}
where $n(\mathbf{r},t)$ is the particle density and $\phi(\mathbf{r},t)$ the phase \cite{pethick2008}.

The Gross-Pitaevskii equation describes the dynamics of the condensate:
\begin{equation}
i \hbar \frac{\partial \psi}{\partial t} = \left( -\frac{\hbar^2}{2m} \nabla^2 + V_{\text{ext}} + g |\psi|^2 \right) \psi,
\end{equation}
where $V_{\text{ext}}$ is the external potential and $g$ is the interaction parameter \cite{nozieres1990}.

\subsection{Multiplicity Theory Overview}
Multiplicity Theory encodes quantum states using prime numbers, enabling the analysis of recursive systems. The interaction matrix $M$ for a prime-encoded system is defined as:
\begin{equation}
M_{ij} = p_i \cdot p_j,
\end{equation}
where $p_i, p_j$ are primes representing distinct quantum states. Eigenvalues $\lambda$ determine system stability:
\begin{equation}
M v = \lambda v,
\end{equation}
where $v$ is the eigenvector \cite{multiplicity2024}.

\section{Helium BEC and Multiplicity Integration}
\subsection{Coherence and Prime-Based Encoding}
The macroscopic wavefunction of Helium BECs is represented as:
\begin{equation}
\psi(t) = \sum_{i=1}^n c_i(t) |p_i\rangle,
\end{equation}
where $c_i(t)$ are time-dependent amplitudes and $|p_i\rangle$ denotes prime-encoded states.

\subsection{Superfluid Dynamics in Recursive Models}
Superfluid vortices in BECs are modeled using recursive dynamics:
\begin{equation}
\psi_{\text{recursive}}(t) = \prod_{i=1}^n \psi^{p_i}_i,
\end{equation}
where the prime powers $p_i$ encode the recursive interactions.

\subsection{Stability Analysis}
The stability of coherence in a Helium BEC is analyzed through the eigenvalues of the interaction matrix:
\begin{equation}
\lambda_i = \prod_{j=1}^m p_j^{\alpha_{ij}},
\end{equation}
where $\alpha_{ij}$ represents the interaction coefficients \cite{pethick2008}.

\section{Applications in Quantum Computing}
Prime-encoded BECs enhance coherence in quantum circuits, enabling robust entanglement and error correction. Quantum gates acting on prime-encoded qubits are represented as:
\begin{equation}
U_p |p_i\rangle = \alpha |p_i\rangle + \beta |p_j\rangle,
\end{equation}
where $\alpha$ and $\beta$ are coefficients determining superposition states \cite{cornell2002}.

\section{Conclusion}
Helium BECs, when integrated with Multiplicity Theory, provide a powerful framework for understanding macroscopic quantum phenomena. The prime-encoded interaction matrices and recursive feedback loops enhance applications in quantum simulations, cryptographic resilience, and astrophysical modeling. Future work will extend this framework to dynamic, multi-scalar systems in high-dimensional spaces.

\bibliographystyle{plain}
\bibliography{references}

\end{document}
